\section{GPU-Native Symbolic Substrate}\label{sec:datalog}

The deterministic Datalog engine supplies CUDA relational operators and device-resident relation storage used across the system. The algorithmic approach---semi-naive fixpoint iteration over stratified programs---is well established; the contribution here is engineering. The ordinary executor is host-orchestrated. For eligible recursive query plans, automatic selection instead uses a certified resident conditional graph: its core runs one recorded launch and synchronization with zero tracked host--device transfers, then returns one pinned terminal receipt carrying status, counters, trace metadata, and schema selection rather than tuple data.

\subsection{Semi-naive evaluation on GPU}

The executor processes an execution plan as strata ordered by the dependency analysis of \Cref{sec:arch}. Non-recursive SCCs evaluate each rule once, dispatching its RIR tree to the appropriate kernel and merging results per head predicate. Recursive SCCs run semi-naive iteration (\Cref{alg:seminaive}): each rule is re-evaluated through delta-rewritten variants (one per recursive scan occurrence, so a self-join such as \texttt{p(X,Y), p(Y,Z)} contributes two variants) and the new tuples are unioned, differenced against the full relation, and deduplicated, all on-device, until no predicate produces new tuples. The profiler records per-operator timing, row counts, and peak memory, feeding the optimizer.

\begin{algorithm}[htbp]
\caption{GPU semi-naive evaluation of a recursive SCC. All buffers are device-resident.}
\label{alg:seminaive}
\begin{algorithmic}[1]
\Require SCC rules $\mathcal{R}$; relation store $S$ (device-resident)
\Ensure least fixpoint of $\mathcal{R}$ merged into $S$
\For{each rule $r \in \mathcal{R}$}
  \State $\Delta_r \gets \textsc{Eval}(r, S) \setminus S$ \Comment{seed $\Delta$; set difference on device}
\EndFor
\Repeat
  \For{each rule $r \in \mathcal{R}$, recursive scan $i$ in $r$}
    \State $T_{r,i} \gets \textsc{Eval}(r[\Delta \text{ at } i], S)$ \Comment{one $\Delta$-variant per self-join}
  \EndFor
  \For{each head predicate $p$}
    \State $\Delta_p \gets \textsc{Dedup}\big(\big(\textstyle\bigcup_{r,i \text{ for } p} T_{r,i}\big) \setminus S_p\big)$
    \State $S_p \gets S_p \cup \Delta_p$ \Comment{union + sorted dedup on device}
  \EndFor
\Until{all $\Delta_p = \emptyset$ \textbf{or} iteration cap reached}
\State free all $\Delta$ buffers
\end{algorithmic}
\end{algorithm}

\subsection{Relational kernels}

xlog implements its relational algebra as direct CUDA kernels rather than wrappers over cuBLAS, cuSPARSE, or Thrust. Bespoke kernels make allocation, synchronization, and host staging explicit and keep relational intermediate buffers on-device. \Cref{tab:kernels} summarizes them.

\begin{table}[htbp]
  \centering
  \caption{Core relational CUDA kernels.}
  \label{tab:kernels}
  \small
  \begin{tabularx}{\linewidth}{@{}l L@{}}
    \toprule
    \textbf{Kernel} & \textbf{Purpose} \\
    \midrule
    \texttt{join}    & Hash join with linked-list collision chains and FNV-1a composite hashing over raw key bytes, so all scalar types share one path; count\,$\to$\,scan\,$\to$\,materialize variants for inner and left-outer joins. \\
    \texttt{sort}    & Stable 4-bit radix sort (8 passes over 32-bit keys): per-pass histogram, prefix sum, scatter. \\
    \texttt{filter}  & Comparison operators for column-vs-constant and column-vs-column, with stream compaction via prefix sum. \\
    \texttt{groupby} & Sorted-input grouped aggregation: boundary detection, then per-group Count/Sum/Min/Max/LogSumExp reduction. \\
    \texttt{dedup}   & Sort-based deduplication and set difference with type-aware equality, including IEEE~754 float handling. \\
    \texttt{set\_ops} & Union (concat\,+\,sort\,+\,dedup) and set difference via sorted-array binary search. \\
    \texttt{scan}    & Blelloch parallel prefix sum, the building block for filter compaction, dedup, and group-by. \\
    \texttt{pack}    & Multi-column key packing into row-major byte arrays with FNV-1a hashing. \\
    \bottomrule
  \end{tabularx}
\end{table}

The filter kernel makes a correctness-critical choice for floating point. Equality uses standard IEEE~754 semantics ($\mathrm{NaN}\neq\mathrm{NaN}$); ordered comparisons use a total-ordering transform that maps an \texttt{f64} bit pattern to an \texttt{i64} whose integer order matches the IEEE~754 \texttt{totalOrder} predicate ($-\mathrm{NaN} < -\mathrm{Inf} < $ negatives $< -0.0 < +0.0 <$ positives $< +\mathrm{Inf} < +\mathrm{NaN}$), so Datalog programs produce deterministic results for all floating-point inputs.

\subsection{Adaptive join planning}

The optimizer performs predicate pushdown using runtime statistics, decomposing filters above joins into left-, right-, and cross-predicates. Lowering chooses join order with a deterministic greedy heuristic. A selectivity pass recognizes triangle and four-cycle shapes, and prior execution statistics can be exported and merged before a later compilation. A general dynamic-programming join ordering is not active; configuration fields that resemble such a threshold do not change this planning path.

\subsection{Worst-case-optimal joins}

Binary-join plans can materialize intermediates asymptotically larger than the final result, the classic pathology for cyclic queries such as triangle enumeration. This is not a corner case for xlog's target workloads: program-analysis queries (points-to, call-graph and reachability inference over code-dependency graphs) can contain cyclic, skewed joins where intermediate size governs cost. xlog addresses this with a worst-case-optimal join (WCOJ) subsystem~\cite{ngo2018wcoj,veldhuizen2014leapfrog,wang2023freejoin}. A pure-Rust hypergraph planner analyzes rule bodies for multiway-join eligibility, derives a deterministic variable ordering from cardinality and selectivity statistics, and emits a \texttt{MultiWayJoin}/\texttt{ChainJoin} node when promotion preserves output semantics; other rules retain their ordinary binary-join plan.

The GPU kernel executes the join via count\,$\to$\,device prefix scan\,$\to$\,materialize over flat sorted-column storage, with a four-byte metadata D2H total and no count-vector transfer. Coverage spans the triangle, 4-cycles, $K$-clique templates ($K{=}5$--$8$), recursive/SCC integration, and helper-relation splitting that exposes inner skew to root-level histograms. A default-on skew classifier decides per query whether to dispatch WCOJ: high-skew inputs take the WCOJ path, while uniform or empty inputs select the binary-join chain. \Cref{sec:wcoj-eval} reports the measured speedup.

\subsection{Runtime optimization}

Interactive and long-running deployments (REPL sessions, iterative solver loops, incremental training) repeatedly evaluate similar queries over relations that change little between calls. A stateless executor rebuilds join indices and recomputes shared subplans each time, paying the dominant cost repeatedly. Three optimizations exploit this between-call stability, each with explicit controls and zero added data-plane transfers. Common-subexpression elimination caches safe deterministic subplans within an evaluation. Adaptive re-optimization adopts a compiler-supplied candidate plan when mis-plan telemetry crosses a configurable ratio, with output-equivalence checks and rollback. And a persistent hash-index manager, keyed on relation generation, schema, and device, retains built indices across repeated sessions with deterministic LRU eviction (\Cref{sec:eval}).

\subsection{Reversible symbols}

Datalog programs operate on strings, but GPU kernels operate on fixed-width numeric columns. xlog interns each unique string to a dense, sequential \texttt{u32} identifier and resolves it back in $O(1)$ at output time. Symbol columns then \emph{are} ordinary \texttt{u32} columns (join, sort, filter, and dedup kernels operate on them with no special-casing) and the string table stays host-side, since variable-length data is needed only at ingestion and output. The \texttt{symbol} type maps to Arrow's \texttt{Dictionary(UInt32, Utf8)}, so export to columnar consumers (cuDF, Polars, DuckDB) is zero-copy while preserving the compact internal representation.
